\begin{lemma}
\label{lemma-localize-flat-dimension-n}
In Situation \ref{situation-flat-dimension-n}.
Let $h : X' \to X$ be an \'etale morphism.
Set $\mathcal{F}' = h^*\mathcal{F}$ and $f' = f \circ h$.
Let $F_n'$ be (\ref{equation-flat-dimension-n})
associated to $(f' : X' \to S, \mathcal{F}')$.
Then $F_n$ is a subfunctor of $F_n'$ and if
$h(X') \supset \text{Ass}_{X/S}(\mathcal{F})$, then $F_n = F'_n$.
\end{lemma}

\begin{proof}
Let $T \to S$ be any morphism. Then $h_T : X'_T \to X_T$ is \'etale as a
base change of the \'etale morphism $g$. For $t \in T$ denote
$Z \subset X_t$ the set of points where $\mathcal{F}_T$ is not
flat over $T$, and similarly denote $Z' \subset X'_t$ the set of
points where $\mathcal{F}'_T$ is not flat over $T$. As
$\mathcal{F}'_T = h_T^*\mathcal{F}_T$ we see that
$Z' = h_t^{-1}(Z)$, see
Morphisms, Lemma \ref{morphisms-lemma-flat-permanence}.
Hence $Z' \to Z$ is an \'etale morphism, so $\dim(Z') \leq \dim(Z)$
(for example by
Descent, Lemma \ref{descent-lemma-dimension-at-point-local}
or just because an \'etale morphism is smooth of relative dimension $0$).
This implies that $F_n \subset F_n'$.

\medskip\noindent
Finally, suppose that $h(X') \supset \text{Ass}_{X/S}(\mathcal{F})$
and that $T \to S$ is a morphism such that $F_n'(T)$ is nonempty, i.e.,
such that $\mathcal{F}'_T$ is flat in dimensions $\geq n$ over $T$.
Pick a point $t \in T$ and let $Z \subset X_t$ and $Z' \subset X'_t$
be as above. To get a contradiction assume that $\dim(Z) \geq n$.
Pick a generic point $\xi \in Z$ corresponding to a component
of dimension $\geq n$. Let $x \in \text{Ass}_{X_t}(\mathcal{F}_t)$
be a generalization of $\xi$. Then $x$ maps to a point of
$\text{Ass}_{X/S}(\mathcal{F})$ by
Divisors, Lemma \ref{divisors-lemma-base-change-relative-assassin} and
Remark \ref{divisors-remark-base-change-relative-assassin}.
Thus we see that $x$ is in the image of $h_T$, say
$x = h_T(x')$ for some $x' \in X'_T$. But $x' \not \in Z'$
as $x \leadsto \xi$ and $\dim(Z') < n$. Hence $\mathcal{F}'_T$
is flat over $T$ at $x'$ which implies that $\mathcal{F}_T$ is flat
at $x$ over $T$ (by
Morphisms, Lemma \ref{morphisms-lemma-flat-permanence}).
Since this holds for every such $x$ we conclude
that $\mathcal{F}_T$ is flat over $T$ at $\xi$ by
Theorem \ref{theorem-check-flatness-at-associated-points}
which is the desired contradiction.
\end{proof}
